\documentclass[11pt]{article}
\usepackage{amsmath,amssymb,booktabs,longtable,geometry,fontspec}
\geometry{margin=1in}
\setmainfont{Times New Roman}[Ligatures=TeX]
\title{Testimony: arguments in propositional form}
\date{Generated by \texttt{lake exe argtex}}
\begin{document}
\maketitle
\noindent Each argument below is a named premise package. Atoms are numbered propositional variables; the legend gives each one's claim, its classification, and the source that carries it. Nothing here asserts that the premises are true --- only that the conclusion does or does not follow from them.

\section*{Sola fide}

Two independent strands, Pauline and dominical. Neither lexical premise carries the argument alone; only their disjunction does.

\subsection*{Reformed (sola fide)}

\begin{longtable}{@{}rp{0.44\textwidth}ll@{}}
\toprule
 & claim & kind & source \\
\midrule
\endhead
$P_{1}$ & Paul's ἔργα νόμου denotes human works in general, not boundary markers & linguistic & \cite[III.xi.19]{calvin-institutes-1960} \\
$P_{2}$ & σῴζω in Luke 7:50 denotes salvation, not physical healing & linguistic & \cite[ad loc. Luke 7:50]{marshall-luke-1978}; Luke 7:47-50 \\
$P_{3}$ & Ephesians 2:8–9 teaches salvation by grace through faith, not of works & textual & Eph 2:8-9; \cite[III.xi.7]{calvin-institutes-1960} \\
$P_{4}$ & Romans 3:28 teaches justification by faith apart from works of the law & textual & Rom 3:28; \cite[III.xi.19]{calvin-institutes-1960} \\
$P_{5}$ & Galatians 2:16 teaches that no one is justified by works of the law & textual & Gal 2:16; \cite[III.xi.19]{calvin-institutes-1960} \\
$P_{6}$ & Romans 4:4–5 contrasts wages owed with a gift reckoned to the one who believes & textual & Rom 4:4-5; \cite[III.xi.18]{calvin-institutes-1960} \\
$P_{7}$ & Titus 3:5 teaches that God saved us not by works done in righteousness & textual & Titus 3:5; \cite[III.xiv.5]{calvin-institutes-1960} \\
$P_{8}$ & Luke 7:50 — Jesus says ‘your faith has saved you’ after declaring sins forgiven & textual & Luke 7:50; \cite[ad loc. Luke 7:50]{marshall-luke-1978} \\
$P_{9}$ & James 2:14–26 targets a barren faith — mere assent — not Paul's doctrine & interpretive & \cite[ad loc. Jas 2:14]{moo-james-2000}; \cite[ad loc. Jas 2:24]{johnson-james-1995}; Jas 2:19 \\
$P_{10}$ & Good works are the fruit and evidence of saving faith, not its ground & theological & \cite[XI.2]{westminster-confession-1647}; \cite[III.xvi.1]{calvin-institutes-1960} \\
$P_{11}$ & Scripture does not contradict itself & theological & \cite[I.vii]{calvin-institutes-1960} \\
$P_{12}$ & Justification is by faith alone & theological & \cite[III.xi.1]{calvin-institutes-1960} \\
$P_{13}$ & James 2:24 is compatible with Paul, using ‘justify’ in a different sense & interpretive & Jas 2:24; \cite[ad loc. Jas 2:24]{moo-james-2000}; \cite[III.xvii.11]{calvin-institutes-1960} \\
$P_{14}$ & Salvation is by grace through faith, and not by works & theological & \cite[III.xi–xviii]{calvin-institutes-1960} \\
\bottomrule
\end{longtable}

\begin{align*}
  \text{(1)} \quad & P_{1} \\
  \text{(2)} \quad & P_{2} \\
  \text{(3)} \quad & P_{3} \\
  \text{(4)} \quad & P_{4} \\
  \text{(5)} \quad & P_{5} \\
  \text{(6)} \quad & P_{6} \\
  \text{(7)} \quad & P_{7} \\
  \text{(8)} \quad & P_{8} \\
  \text{(9)} \quad & P_{9} \\
  \text{(10)} \quad & P_{10} \\
  \text{(11)} \quad & P_{11} \\
  \text{(12)} \quad & (P_{4} \land P_{5} \land P_{1}) \rightarrow P_{12} \\
  \text{(13)} \quad & (P_{8} \land P_{2}) \rightarrow P_{12} \\
  \text{(14)} \quad & (P_{9} \land P_{10}) \rightarrow P_{13} \\
  \text{(15)} \quad & (P_{12} \land P_{3} \land P_{6} \land P_{7} \land P_{13} \land P_{11}) \rightarrow P_{14} \\
  \midrule[0pt] \vdash \quad & P_{14}
\end{align*}

\emph{No premise rests on scripture alone.}
\subsection*{New Perspective on Paul}

\begin{longtable}{@{}rp{0.44\textwidth}ll@{}}
\toprule
 & claim & kind & source \\
\midrule
\endhead
$P_{1}$ & Paul's ἔργα νόμου denotes Jewish covenant boundary markers, not works in general & linguistic & \cite{dunn-new-perspective-2005}; \cite{wright-what-paul-said-1997} \\
$P_{2}$ & σῴζω in Luke 7:50 denotes salvation, not physical healing & linguistic & \cite[ad loc. Luke 7:50]{marshall-luke-1978}; Luke 7:47-50 \\
$P_{3}$ & Ephesians 2:8–9 teaches salvation by grace through faith, not of works & textual & Eph 2:8-9; \cite[III.xi.7]{calvin-institutes-1960} \\
$P_{4}$ & Romans 3:28 teaches justification by faith apart from works of the law & textual & Rom 3:28; \cite[III.xi.19]{calvin-institutes-1960} \\
$P_{5}$ & Galatians 2:16 teaches that no one is justified by works of the law & textual & Gal 2:16; \cite[III.xi.19]{calvin-institutes-1960} \\
$P_{6}$ & Romans 4:4–5 contrasts wages owed with a gift reckoned to the one who believes & textual & Rom 4:4-5; \cite[III.xi.18]{calvin-institutes-1960} \\
$P_{7}$ & Titus 3:5 teaches that God saved us not by works done in righteousness & textual & Titus 3:5; \cite[III.xiv.5]{calvin-institutes-1960} \\
$P_{8}$ & Luke 7:50 — Jesus says ‘your faith has saved you’ after declaring sins forgiven & textual & Luke 7:50; \cite[ad loc. Luke 7:50]{marshall-luke-1978} \\
$P_{9}$ & James 2:14–26 targets a barren faith — mere assent — not Paul's doctrine & interpretive & \cite[ad loc. Jas 2:14]{moo-james-2000}; \cite[ad loc. Jas 2:24]{johnson-james-1995}; Jas 2:19 \\
$P_{10}$ & Good works are the fruit and evidence of saving faith, not its ground & theological & \cite[XI.2]{westminster-confession-1647}; \cite[III.xvi.1]{calvin-institutes-1960} \\
$P_{11}$ & Scripture does not contradict itself & theological & \cite[I.vii]{calvin-institutes-1960} \\
$P_{12}$ & Justification is by faith alone & theological & \cite[III.xi.1]{calvin-institutes-1960} \\
$P_{13}$ & James 2:24 is compatible with Paul, using ‘justify’ in a different sense & interpretive & Jas 2:24; \cite[ad loc. Jas 2:24]{moo-james-2000}; \cite[III.xvii.11]{calvin-institutes-1960} \\
$P_{14}$ & Salvation is by grace through faith, and not by works & theological & \cite[III.xi–xviii]{calvin-institutes-1960} \\
\bottomrule
\end{longtable}

\begin{align*}
  \text{(1)} \quad & \lnot P_{1} \\
  \text{(2)} \quad & P_{2} \\
  \text{(3)} \quad & P_{3} \\
  \text{(4)} \quad & P_{4} \\
  \text{(5)} \quad & P_{5} \\
  \text{(6)} \quad & P_{6} \\
  \text{(7)} \quad & P_{7} \\
  \text{(8)} \quad & P_{8} \\
  \text{(9)} \quad & P_{9} \\
  \text{(10)} \quad & P_{10} \\
  \text{(11)} \quad & P_{11} \\
  \text{(12)} \quad & (P_{4} \land P_{5} \land P_{1}) \rightarrow P_{12} \\
  \text{(13)} \quad & (P_{8} \land P_{2}) \rightarrow P_{12} \\
  \text{(14)} \quad & (P_{9} \land P_{10}) \rightarrow P_{13} \\
  \text{(15)} \quad & (P_{12} \land P_{3} \land P_{6} \land P_{7} \land P_{13} \land P_{11}) \rightarrow P_{14} \\
  \midrule[0pt] \vdash \quad & P_{14}
\end{align*}

\emph{No premise rests on scripture alone.}
\subsection*{Tridentine (Council of Trent, Session VI)}

\begin{longtable}{@{}rp{0.44\textwidth}ll@{}}
\toprule
 & claim & kind & source \\
\midrule
\endhead
$P_{1}$ & Ephesians 2:8–9 teaches salvation by grace through faith, not of works & textual & Eph 2:8-9; \cite[III.xi.7]{calvin-institutes-1960} \\
$P_{2}$ & Romans 3:28 teaches justification by faith apart from works of the law & textual & Rom 3:28; \cite[III.xi.19]{calvin-institutes-1960} \\
$P_{3}$ & Galatians 2:16 teaches that no one is justified by works of the law & textual & Gal 2:16; \cite[III.xi.19]{calvin-institutes-1960} \\
$P_{4}$ & Titus 3:5 teaches that God saved us not by works done in righteousness & textual & Titus 3:5; \cite[III.xiv.5]{calvin-institutes-1960} \\
$P_{5}$ & Scripture does not contradict itself & theological & \cite[I.vii]{calvin-institutes-1960} \\
$P_{6}$ & Works performed in grace merit an increase of justification & theological & \cite[Trent, Session VI (1547), Decree on Justification, ch. 16]{tanner-decrees-1990} \\
$P_{7}$ & Salvation is by grace through faith, and not by works & theological & \cite[III.xi–xviii]{calvin-institutes-1960} \\
\bottomrule
\end{longtable}

\begin{align*}
  \text{(1)} \quad & P_{1} \\
  \text{(2)} \quad & P_{2} \\
  \text{(3)} \quad & P_{3} \\
  \text{(4)} \quad & P_{4} \\
  \text{(5)} \quad & P_{5} \\
  \text{(6)} \quad & P_{6} \\
  \text{(7)} \quad & P_{6} \rightarrow \lnot P_{7} \\
  \midrule[0pt] \vdash \quad & P_{7}
\end{align*}

\emph{No premise rests on scripture alone.}
\section*{Sola scriptura}

Seeded. The self-refutation objection is encoded alongside the position it tells against.

\subsection*{Protestant (sola scriptura)}

\begin{longtable}{@{}rp{0.44\textwidth}ll@{}}
\toprule
 & claim & kind & source \\
\midrule
\endhead
$P_{1}$ & 2 Timothy 3:16 — all scripture is God-breathed and profitable & textual & 2 Tim 3:16; \cite[I.vii.1]{calvin-institutes-1960} \\
$P_{2}$ & 2 Timothy 3:17 — that the man of God may be complete, equipped for every good work & textual & 2 Tim 3:17; \cite[I.vii.1]{calvin-institutes-1960} \\
$P_{3}$ & Mark 7:8–13 — Jesus rebukes tradition that nullifies God's command & textual & Mark 7:8-13; \cite[IV.x.8]{calvin-institutes-1960} \\
$P_{4}$ & Acts 17:11 — the Bereans tested apostolic preaching against scripture & textual & Acts 17:11; \cite[I.vii.2]{calvin-institutes-1960} \\
$P_{5}$ & Scripture is a sufficient rule of faith & theological & \cite[I.vii]{calvin-institutes-1960} \\
$P_{6}$ & Scripture is clear on what is necessary for salvation & theological & \cite[I.vii.5]{calvin-institutes-1960} \\
$P_{7}$ & Scripture itself teaches that scripture is the sole infallible rule & interpretive & \cite[I.vii.4]{calvin-institutes-1960} \\
$P_{8}$ & Scripture is the sole infallible rule of faith & theological & \cite[I.vii–ix]{calvin-institutes-1960} \\
\bottomrule
\end{longtable}

\begin{align*}
  \text{(1)} \quad & P_{1} \\
  \text{(2)} \quad & P_{2} \\
  \text{(3)} \quad & P_{3} \\
  \text{(4)} \quad & P_{4} \\
  \text{(5)} \quad & P_{5} \\
  \text{(6)} \quad & P_{6} \\
  \text{(7)} \quad & P_{7} \\
  \text{(8)} \quad & (P_{1} \land P_{2} \land P_{5} \land P_{6} \land P_{7}) \rightarrow P_{8} \\
  \midrule[0pt] \vdash \quad & P_{8}
\end{align*}

\emph{No premise rests on scripture alone.}
\subsection*{Catholic/Orthodox (scripture with tradition)}

\begin{longtable}{@{}rp{0.44\textwidth}ll@{}}
\toprule
 & claim & kind & source \\
\midrule
\endhead
$P_{1}$ & 2 Timothy 3:16 — all scripture is God-breathed and profitable & textual & 2 Tim 3:16; \cite[I.vii.1]{calvin-institutes-1960} \\
$P_{2}$ & 2 Thessalonians 2:15 — hold to the traditions taught by word or letter & textual & 2 Thess 2:15; \cite[Trent, Session IV (1546), Decree on Sacred Books and Traditions]{tanner-decrees-1990} \\
$P_{3}$ & The magisterium is an infallible interpreter of scripture & theological & \cite[Trent, Session IV (1546), Decree on the Vulgate Edition]{tanner-decrees-1990} \\
$P_{4}$ & Scripture is the sole infallible rule of faith & theological & \cite[I.vii–ix]{calvin-institutes-1960} \\
\bottomrule
\end{longtable}

\begin{align*}
  \text{(1)} \quad & P_{1} \\
  \text{(2)} \quad & P_{2} \\
  \text{(3)} \quad & P_{3} \\
  \text{(4)} \quad & (P_{2} \land P_{3}) \rightarrow \lnot P_{4} \\
  \midrule[0pt] \vdash \quad & P_{4}
\end{align*}

\emph{No premise rests on scripture alone.}
\subsection*{Self-refutation objection to sola scriptura}

\begin{longtable}{@{}rp{0.44\textwidth}ll@{}}
\toprule
 & claim & kind & source \\
\midrule
\endhead
$P_{1}$ & A doctrine is binding only if scripture teaches it & theological & \cite[IV.x.8]{calvin-institutes-1960} \\
$P_{2}$ & Scripture itself teaches that scripture is the sole infallible rule & interpretive & \cite[I.vii.4]{calvin-institutes-1960} \\
$P_{3}$ & Scripture is the sole infallible rule of faith & theological & \cite[I.vii–ix]{calvin-institutes-1960} \\
\bottomrule
\end{longtable}

\begin{align*}
  \text{(1)} \quad & P_{1} \\
  \text{(2)} \quad & \lnot P_{2} \\
  \text{(3)} \quad & (P_{1} \land \lnot P_{2}) \rightarrow \lnot P_{3} \\
  \midrule[0pt] \vdash \quad & \lnot P_{3}
\end{align*}

\emph{No premise rests on scripture alone.}
\section*{Born in Bethlehem (Micah 5:2)}

Matthew 2:5--6 quotes Micah 5:2 as grounds for the Messiah's birthplace.

\subsection*{Christian predictive reading of Micah 5:2}

\begin{longtable}{@{}rp{0.44\textwidth}ll@{}}
\toprule
 & claim & kind & source \\
\midrule
\endhead
$P_{1}$ & Micah 5:2 is a forward-looking Messianic prediction & interpretive & \cite[ad loc. Mic 5:2]{keil-delitzsch-minor-prophets-1949} \\
$P_{2}$ & Matthew 2:6 quotes Micah 5:2 & textual & \cite[app. Matt 2:6]{na28-2012}; \cite{ubs5-2014} \\
$P_{3}$ & Matthew's quotation intends predictive fulfilment & interpretive & \cite[ad loc. Matt 2:6]{france-matthew-2007} \\
$P_{4}$ & Jesus of Nazareth was born in Bethlehem & historical & Matt 2:1; Luke 2:4-7 \\
$P_{5}$ & The Messiah must be born in Bethlehem & interpretive & \cite[ad loc. Mic 5:2]{keil-delitzsch-minor-prophets-1949}; \cite[ad loc. Matt 2:6]{france-matthew-2007} \\
$P_{6}$ & Jesus of Nazareth satisfies the Bethlehem criterion & interpretive & \cite[ad loc. Matt 2:6]{france-matthew-2007} \\
\bottomrule
\end{longtable}

\begin{align*}
  \text{(1)} \quad & P_{1} \\
  \text{(2)} \quad & P_{2} \\
  \text{(3)} \quad & P_{3} \\
  \text{(4)} \quad & P_{4} \\
  \text{(5)} \quad & (P_{1} \land P_{3}) \rightarrow P_{5} \\
  \text{(6)} \quad & (P_{5} \land P_{4}) \rightarrow P_{6} \\
  \midrule[0pt] \vdash \quad & P_{6}
\end{align*}

\emph{Grounded in scripture alone:} Jesus of Nazareth was born in Bethlehem.
\subsection*{Critical reading of Micah 5:2 as a near-term oracle}

\begin{longtable}{@{}rp{0.44\textwidth}ll@{}}
\toprule
 & claim & kind & source \\
\midrule
\endhead
$P_{1}$ & Matthew 2:6 quotes Micah 5:2 & textual & \cite[app. Matt 2:6]{na28-2012}; \cite{ubs5-2014} \\
$P_{2}$ & Micah 5:2 is a near-term oracle about a contemporary Judaean ruler & interpretive & \cite{brown-birth-messiah-1993} \\
$P_{3}$ & Micah 5:2 is a forward-looking Messianic prediction & interpretive & \cite[ad loc. Mic 5:2]{keil-delitzsch-minor-prophets-1949} \\
$P_{4}$ & Matthew's quotation intends predictive fulfilment & interpretive & \cite[ad loc. Matt 2:6]{france-matthew-2007} \\
$P_{5}$ & The Messiah must be born in Bethlehem & interpretive & \cite[ad loc. Mic 5:2]{keil-delitzsch-minor-prophets-1949}; \cite[ad loc. Matt 2:6]{france-matthew-2007} \\
$P_{6}$ & Jesus of Nazareth was born in Bethlehem & historical & Matt 2:1; Luke 2:4-7 \\
$P_{7}$ & Jesus of Nazareth satisfies the Bethlehem criterion & interpretive & \cite[ad loc. Matt 2:6]{france-matthew-2007} \\
\bottomrule
\end{longtable}

\begin{align*}
  \text{(1)} \quad & P_{1} \\
  \text{(2)} \quad & P_{2} \\
  \text{(3)} \quad & \lnot P_{3} \\
  \text{(4)} \quad & (P_{3} \land P_{4}) \rightarrow P_{5} \\
  \text{(5)} \quad & (P_{5} \land P_{6}) \rightarrow P_{7} \\
  \midrule[0pt] \vdash \quad & P_{7}
\end{align*}

\emph{Grounded in scripture alone:} Jesus of Nazareth was born in Bethlehem.
\section*{Born of a virgin (Isaiah 7:14)}

A single-stranded argument, and so the weaker one: defeating the lexical premise defeats it outright.

\subsection*{Christian predictive reading of Isaiah 7:14}

\begin{longtable}{@{}rp{0.44\textwidth}ll@{}}
\toprule
 & claim & kind & source \\
\midrule
\endhead
$P_{1}$ & Isaiah 7:14 is a Messianic prediction of a virgin birth & interpretive & \cite[ad loc. Isa 7:14]{motyer-isaiah-1993} \\
$P_{2}$ & עַלְמָה in Isaiah 7:14 denotes a virgin, not merely a young woman & linguistic & \cite[ad loc. Isa 7:14]{motyer-isaiah-1993}; \cite[app. Isa 7:14]{bhs-1997} \\
$P_{3}$ & The Septuagint renders עַלְמָה at Isaiah 7:14 as παρθένος & textual & \cite[app. Matt 1:23]{na28-2012}; \cite[ad loc. Matt 1:23]{france-matthew-2007} \\
$P_{4}$ & Matthew 1:23 quotes Isaiah 7:14 & textual & \cite[app. Matt 1:23]{na28-2012}; \cite{ubs5-2014} \\
$P_{5}$ & Matthew's quotation intends the virgin conception as fulfilment & interpretive & \cite[ad loc. Matt 1:23]{france-matthew-2007} \\
$P_{6}$ & Mary conceived Jesus while a virgin & historical & Matt 1:18-25; Luke 1:26-38 \\
$P_{7}$ & Matthew and Luke are independent traditions agreeing on the virgin conception & historical & \cite[26--38]{brown-birth-messiah-1993} \\
$P_{8}$ & The Messiah must be born of a virgin & interpretive & \cite[ad loc. Isa 7:14]{motyer-isaiah-1993}; \cite[ad loc. Matt 1:23]{france-matthew-2007} \\
$P_{9}$ & Jesus of Nazareth satisfies the virgin-birth criterion & interpretive & \cite[ad loc. Matt 1:23]{france-matthew-2007} \\
\bottomrule
\end{longtable}

\begin{align*}
  \text{(1)} \quad & P_{1} \\
  \text{(2)} \quad & P_{2} \\
  \text{(3)} \quad & P_{3} \\
  \text{(4)} \quad & P_{4} \\
  \text{(5)} \quad & P_{5} \\
  \text{(6)} \quad & P_{6} \\
  \text{(7)} \quad & P_{7} \\
  \text{(8)} \quad & (P_{1} \land P_{2} \land P_{5}) \rightarrow P_{8} \\
  \text{(9)} \quad & (P_{8} \land P_{6}) \rightarrow P_{9} \\
  \midrule[0pt] \vdash \quad & P_{9}
\end{align*}

\emph{Grounded in scripture alone:} Mary conceived Jesus while a virgin.
\subsection*{Critical reading of Isaiah 7:14 as a near-term sign}

\begin{longtable}{@{}rp{0.44\textwidth}ll@{}}
\toprule
 & claim & kind & source \\
\midrule
\endhead
$P_{1}$ & The Septuagint renders עַלְמָה at Isaiah 7:14 as παρθένος & textual & \cite[app. Matt 1:23]{na28-2012}; \cite[ad loc. Matt 1:23]{france-matthew-2007} \\
$P_{2}$ & Matthew 1:23 quotes Isaiah 7:14 & textual & \cite[app. Matt 1:23]{na28-2012}; \cite{ubs5-2014} \\
$P_{3}$ & Isaiah 7:14 is a near-term sign to Ahaz, fulfilled in Isaiah's generation & interpretive & \cite{brown-birth-messiah-1993} \\
$P_{4}$ & עַלְמָה in Isaiah 7:14 denotes a virgin, not merely a young woman & linguistic & \cite[ad loc. Isa 7:14]{motyer-isaiah-1993}; \cite[app. Isa 7:14]{bhs-1997} \\
$P_{5}$ & Isaiah 7:14 is a Messianic prediction of a virgin birth & interpretive & \cite[ad loc. Isa 7:14]{motyer-isaiah-1993} \\
$P_{6}$ & Matthew's quotation intends the virgin conception as fulfilment & interpretive & \cite[ad loc. Matt 1:23]{france-matthew-2007} \\
$P_{7}$ & The Messiah must be born of a virgin & interpretive & \cite[ad loc. Isa 7:14]{motyer-isaiah-1993}; \cite[ad loc. Matt 1:23]{france-matthew-2007} \\
$P_{8}$ & Mary conceived Jesus while a virgin & historical & Matt 1:18-25; Luke 1:26-38 \\
$P_{9}$ & Jesus of Nazareth satisfies the virgin-birth criterion & interpretive & \cite[ad loc. Matt 1:23]{france-matthew-2007} \\
\bottomrule
\end{longtable}

\begin{align*}
  \text{(1)} \quad & P_{1} \\
  \text{(2)} \quad & P_{2} \\
  \text{(3)} \quad & P_{3} \\
  \text{(4)} \quad & \lnot P_{4} \\
  \text{(5)} \quad & \lnot P_{5} \\
  \text{(6)} \quad & (P_{5} \land P_{4} \land P_{6}) \rightarrow P_{7} \\
  \text{(7)} \quad & (P_{7} \land P_{8}) \rightarrow P_{9} \\
  \midrule[0pt] \vdash \quad & P_{9}
\end{align*}

\emph{Grounded in scripture alone:} Mary conceived Jesus while a virgin.

\bibliographystyle{plain}
\bibliography{references}
\end{document}
